{\Large\textbf{2\quad Black box}}

{\large\textbf{2.1\quad Introduction}}

You have a rigid mechanical black box consisting of a container of mass $m_1$.
Inside the container there is a load of mass $m_2$ that hangs on an effectively
massless spring of stiffness $k_1$ from the ceiling of the box. Another mass
$m_3$ is hanged to the mass $m_2$ via another massless spring of stiffness
$k_2$. There is a small viscous drag which depends on the velocity of the
objects. The gravity of Earth is $g = 9.81\,\mathrm{m/s^2}$ and is parallel to
the sides of the box.

\begin{center}\includegraphics[width=0.40\linewidth]{figures/EuPhO_2020_Q5_fig1.png}\end{center}

The box can be moved up or down with a piece-wise constant acceleration. The
acceleration pattern can be programmed through input by giving the duration (in
seconds) and acceleration (in $\mathrm{m/s^2}$) for each step. The simulation
shows in ”real time” the force $F$ exerted to the box that is needed to
maintain the given acceleration at the current moment of time, together with
the reading of time. The simulation will also output the readings to a text
file in the same folder as the program. All simulations will always start with
the same initial configuration for the masses.

\textit{Note}: Every measurement of force $F$ has a small random error. The
springs are linear for reasonably small deformations, but nonlinear for larger
deformations. The values $k_1$ and $k_2$ are defined to be the stiffness of
each spring for small deformations near equilibrium when the box is at rest.
Force $F$ and acceleration are considered to be positive if they are directed
upwards. The side length of the box is $0.6\,\mathrm{m}$ and the box is
initially in the middle of a room with height $3\,\mathrm{m}$. An experiment
ends automatically if the box hits the ceiling or the floor, or if any of the
masses collide with the box or with the other mass. The figure is not drawn to
scale.

{\large\textbf{2.2\quad Task}}

The task is to determine all the parameters: $m_1$, $m_2$, $m_3$, $k_1$, $k_2$.
You do not need to provide an error analysis for these results.

As with all experiments, you must provide clearly labelled tables of data,
clearly labelled graphs, and sufficient formulae derivations to make it clear
what you have measured, and how you are deriving your results.

{\large\textbf{2.3\quad Program Interface}}

Initially, the program asks for a sequence of input actions. You have the
following possibilities.
\begin{itemize}
\item Enter two numbers and press \textbf{return} to add a step to the
acceleration pattern, for example: \texttt{1.5 -0.4}\\
The first number should be the \textbf{duration} of the step in seconds (must
be a multiple of $0.01\,\mathrm{s}$) and the second number should be the
\textbf{acceleration} in $\mathrm{m/s^2}$ (must be between $-30$ and $30$).
\item Enter \texttt{repeat} and an integer and press \textbf{return} to repeat
actions, for example: \texttt{repeat 10}\\
The integer should be the \textbf{number of times} you want to repeat actions.
Every repeat action should end with an \texttt{endrepeat} action (see below).
\item Enter \texttt{endrepeat} to end repeating actions. If you start the
experiment, all actions between \texttt{repeat} and \texttt{endrepeat} will be
repeated a given number of times. You cannot repeat actions inside another
repeat.
\item Enter \texttt{sample} and a number and press \textbf{return} to change
the sampling time, for example: \texttt{sample 0.4}\\
The number should be the new \textbf{sampling time} which is the time after
which every new reading is output to the text file. The sampling time must be a
multiple of $0.01\,\mathrm{s}$, which is also the default sampling time.
\item Enter \texttt{begin} to finish the sequence and start the experiment.
\end{itemize}

You can also write multiple actions on the same line and then press
\textbf{return}. For example, you can enter

\texttt{sample 0.4 repeat 10 1.5 0.4 1.5 -0.4 endrepeat begin}

to start an experiment where you change the sampling time to $0.4\,\mathrm{s}$
and accelerate the box respectively with $a = 0.4\,\mathrm{m/s^2}$ and
$a = -0.4\,\mathrm{m/s^2}$ ten times.

If you enter an invalid input, you will get one of the following error messages
and you can try to enter an action again.
\begin{itemize}
\item If acceleration is out of range:\\
\texttt{Acceleration is out of range.}
\item If duration of acceleration is out of range:\\
\texttt{Duration is out of range.}
\item If sampling time is out of range:\\
\texttt{Sampling time is out of range.}
\item If the number of repeat times is out of range:\\
\texttt{Number of repeat times is out of range.}
\item If you try to repeat actions inside another repeat action:\\
\texttt{Cannot repeat actions inside another repeat.}
\item If you try to end repeat without a repeat action to end:\\
\texttt{Cannot end repeat outside repeat.}
\item In all other cases:\\
\texttt{Invalid entry.}
\end{itemize}

After you enter \texttt{begin}, the program will ask you for a name for the
output file with the prompt

\texttt{Enter name for output file (e.g. "results"). You should use Latin
letters and numbers because some special characters are not allowed.}

Enter a name and press \textbf{return}. You are advised to use only Latin
letters and numbers for the name. Other characters may or may not be allowed in
the filename and in case of an invalid filename the readings will not be saved.
The readings will be saved in a \texttt{.txt} file with the given name in the
same folder as the program.

After this, the program will display

\texttt{Begin experiment.}

and start the experiment. The program will then display the current time since
the beginning of the experiment (\texttt{Time (s)}), measured value of force
$F$ (\texttt{Force (N)}) and acceleration of the box
(\texttt{Accel (m/s\textasciicircum 2)}). The readings will be similarly
displayed in the text file.

The program will then display one of the following messages.
\begin{itemize}
\item If the experiment ended successfully:\\
\texttt{Experiment ended successfully.}
\item If the box hit the ceiling:\\
\texttt{The box hit the ceiling. Experiment ended.}
\item If the box hit the floor:\\
\texttt{The box hit the floor. Experiment ended.}
\item If the masses inside the box collided or one of the masses inside the box
collided with the box:\\
\texttt{Masses and/or the box collided. Experiment ended.}
\end{itemize}

After the experiment ends, you can start another experiment.
